%%%%%%%%%%%%%%%%%%%%%%%%%%%%%%%%%%%%%%%%%%%%%%%%%%
\subsection{Problem Formulation}\label{remdo:sec:formulation}
The standard optimization problem is expressed \citep{martins_engineering_2022,papalambros_principles_2017}:
\begin{equation}
\label{remdo:eq:standard-opt-prob}
\begin{aligned}
    \text{min}~~& \vec{J}(\vec{x},~\vec{p}) \\
    \text{by varying}~~&\vec{x} \\
    \text{subject to}~~ &\vec{x}_{LB} \leq \vec{x} \leq \vec{x}_{UB} \\
    &\vec{g}(\vec{x},~\vec{p}) \geq 0\\
    &\vec{h}(\vec{x},~\vec{p})= 0\\
\end{aligned}
\end{equation}
An optimization algorithm adjusts the design variable vector $\vec{x}$ within lower and upper bounds $\vec{x}_{LB}$, $\vec{x}_{UB}$ to minimize the objective(s) $\vec{J}$, subject to inequality and equality constraints $\vec{g}$ and $\vec{h}$.
Parameters $\vec{p}$ are constant inputs required to compute the objective or constraints.

%%%%%%%%%%%%%%%%%%%%%%%%
\subsubsection{Objective Function}\label{remdo:sec:obj}
For the single-objective formulation, the optimization minimizes the levelized cost of energy ($LCOE$), the metric most commonly used by WEC developers to evaluate market viability \citep{trueworthy_wave_2020} and computed by the economics module \citep{mccabe_development_2026}.

\ifdefined\DISSERTATION
    LCOE is adopted for its simplicity, although this choice brings three major limitations.
    First, because LCOE balances the device's lifetime cost with its lifetime energy production, it only measures economic viability if each Joule generated has equal value.
    On real grids, the value of electricity varies temporally and spatially, especially with growing electrification \citep{akdemir_opportunities_2023,bhattacharya_timing_2021,pennock_temporal_2022},
    so capturing nonuniform sources of WEC value such as consistency or complementarity would require modifications (e.g.\ those proposed by \citet{mccabe_system_2023}) outside the scope of the present model.
    Second, LCOE has high uncertainty for early-stage designs: the economic model assumes only structural and PTO costs scale with design \citep{mccabe_development_2026}, so changes to fixed cost components or the fixed charge rate ($FCR$) could influence the optimal design.
    Third, LCOE has limited relevance for the off-grid and micro-grid use cases such as offshore aquaculture, autonomous underwater vehicle charging, ocean data collection, that have grown popular since the 2019 launch of the DOE's ``Powering the Blue Economy'' initiative \citep{livecchi_powering_2019}.
    For a fixed power requirement, only cost reduction (not power increase) drives value, and some applications such as ocean observation are not financed at all, making the fixed charge rate obsolete \citep{jenne_powering_2021}.
    
    In the multi-objective formulation, instead of $LCOE$, the two objectives are average electrical power $P_{\text{avg,elec}}$ and design-dependent capital cost $C_{\text{design}}$.
    This directly addresses the second and third LCOE limitations: the fixed charge rate and inaccurately-scaled cost components can no longer influence the optimum, and the minimum-capital design for a given power requirement is easily identified.
\else
    LCOE is adopted for its simplicity, though it carries three limitations:
    it values every Joule equally, ignoring the temporal and spatial variation of electricity value on real grids \citep{akdemir_opportunities_2023,mccabe_system_2023};
    it has high early-stage uncertainty, since costs other than structural and PTO costs are difficult to predict early on and thus modeled as fixed with respect to design \citep{mccabe_development_2026};
    and it is less relevant for the off-grid \emph{Powering the Blue Economy} applications \citep{livecchi_powering_2019,jenne_powering_2021} where power requirements are fixed or projects are unfinanced.
    The multi-objective formulation addresses the latter two limitations by replacing $LCOE$ with average electrical power $P_{\text{avg,elec}}$ and design-dependent capital cost $C_{\text{design}}$.
    This removes the influence of the fixed charge rate and inaccurately-scaled costs and exposes the minimum-capital design for a given power requirement.
\fi

The objective function is expressed:
\begin{equation}
\label{remdo:eq:obj-fcn}
\vec{J}(\vec{x}, \vec{p}) = \begin{cases} LCOE\ (\vec{x}, \vec{p}), & \text{Single objective} \\
[C_{\text{design}} (\vec{x}, \vec{p}),\ P_{\text{avg,elec}} (\vec{x}, \vec{p})]^T, & \text{Multi-objective}
\end{cases}
\end{equation}
%The objectives of this optimization problem are to minimize the LCOE and the power coefficient of variation of the RM3 WEC. LCOE, the cost per unit energy over the lifetime of the device, was chosen as an objective function because it is often the metric used to assess and compare energy production methods. A low LCOE indicates the project is competitive with conventional energy production processes and better able to generate affordable energy for consumers. The second objective, power coefficient of variation, is a measure of the fluctuation in power output. Minimizing power variation provides more consistent output and reduces the amount of energy storage required.

%The coefficient of variation of power $c_v$ normalizes the the standard deviation of power across different sea states $\sigma$ by the average power across those sea states $\mu$:
% \begin{equation}
%     c_v = \frac{\sigma}{\mu}.
%     \label{remdo:cv}
% \end{equation}

%%%%%%%%%%%%%%%%%%%%%%%%
\subsubsection{Requirements}\label{remdo:sec:requirements}
The requirements that the optimization must enforce significantly shape the formulation.
Here, requirements include both traditional design requirements (model-agnostic criteria ensuring the design meets stakeholder needs and is physically realizable) and criteria that facilitate the optimization itself (such as keeping intermediate designs within the model's range of validity).
Constraints are the most common but not the only enforcement mechanism.
\ifdefined\DISSERTATION
    Explicitly outlining the requirements and resulting formulation logic serves three purposes:
    (1) to ground the optimization within the engineering design process, where requirements drive decisions;
    (2) to clarify a part of the optimization process that is critical to result quality yet often omitted;
    and (3) to illustrate how the MDO mindset systematically yields a ``clever'' formulation that overcomes common difficulties in large-scale optimization.
\else
    Outlining the requirements and the resulting formulation logic grounds the optimization within the engineering design process, clarifies a step that is often omitted despite its necessity, and illustrates how the MDO mindset yields a formulation that overcomes common large-scale optimization difficulties.
\fi

\Cref{remdo:tab:requirements} lists each requirement with its priority and enforcement mechanism.
\citet{mccabe_development_2026} develops the physical and modeling basis for each requirement, such as the underlying geometry, dynamics, hydrostatics, structural, and economic relations.
\newcounter{reqnumber}
\singleColMacro{
    \begin{longtable}{cP{0.25\linewidth}P{0.27\linewidth}P{0.08\linewidth}L{0.28\linewidth}}
          \textbf{\#}&\textbf{Requirement}&  \textbf{Description}& \textbf{Priority}&\textbf{Enforcement}\\
         \hline
  \requirement{float_spar_diam}&$D_s < D_{f,in}$&  Float spar diameter& 1&Parameter $D_{f,in}/D_s$\\
  \requirement{float_in_out_diam}&$D_{f,in}<D_{f}$& Float outer/inner diameter& 1&Lin. ineq. constraint $g_{L,2}$\\
  \requirement{float_spar_draft}&$T_{f,2}<T_s$&  Float spar draft& 1&Lin. ineq. constraint $g_{L,3}$\\
  \requirement{float_seafloor_overflow}&$h-T_s>\Delta z_{s,\text{min}}$&  Float seafloor + numerical overflow& 1&Lin. ineq. constraint $g_{L,5}$\\
  \requirement{spar_seafloor_overflow}&$h-T_{f,2}>\Delta z_{f,\text{min}}$&  Spar seafloor + numerical overflow& 1&Lin. ineq. constraint $g_{L,6}$\\
  \hline
  \requirement{float_base_diameter}& $D_{f,in} <D_{f,b} < D_f$& Float base diameter& 2&Parameter $D_{f,b}/D_f$ \\
  \requirement{float_draft}& $T_{f,1}<T_{f,2}$ & Float draft & 2 & Parameter $T_{f,1}/T_{f,2}$ \\
  \requirement{float_submergence}& $T_{f,2} < h_f$& Float submergence& 2&Parameter $T_{f,2}/h_f$\\
  \requirement{spar_damping_ratio}& $\zeta > \zeta_{\text{min}}$& Spar damping ratio& 2&Parameters $D_d/D_s$, $T_s/D_s$, $h_d/D_s$\\
  \requirement{spar_similarity}&$\omega_n<\omega_{n,\text{max}}$& Spar similarity& 2&Lin. ineq. constraint $g_{L,1}$\\
  \requirement{positive_net_power}&$P_{\text{avg}}>0$& Positive net power& 2&Nonlin. ineq. constraint $g_{NL,14}$\\
  \hline
  \requirement{clearance_static}& $h_{fs,\text{clear}}>0$& Float spar tube clearance (static)& 3&Design variable bound\\
  \requirement{tops_static}&$h_{fs,\text{up}}>0$& Float spar tops (static)& 3&Lin. ineq. constraint $g_{L,4}$\\
  \requirement{base_static}& $h_{fs,\text{down}}>0$& Float spar base (static)& 3&Lin. ineq. constraint $g_{L,3}$\\
  \requirement{damping_plate_thickness}& $t_d < h_d$& Maximum damping plate thickness & 3 & Lin. ineq. constraint $g_{L,7}$\\
  \requirement{float_stiffener_height}& $h_{\text{stiff},f} < h_f / 2$& Maximum float stiffener height& 3 &Lin. ineq. constraint $g_{L,8}$\\
  \requirement{tops_dynamic}&$h_{fs,\text{up}}>|\xi_f-\xi_s|$& Float spar tops (dynamics)& 3&Nonlin. ineq. constraint $g_{NL,19}$\\
  \requirement{base_dynamic}& $h_{fs,\text{down}}>|\xi_f-\xi_s|$& Float spar base (dynamic)& 3&Nonlin. ineq. constraint $g_{NL,20}$\\
  \requirement{clearance_dynamic}& $h_{fs,\text{clear}}>|\xi_f-\xi_s|$& Float spar tube clearance (dynamic)& 3&Nonlin. ineq. constraint $g_{NL,21}$\\
  \requirement{linear_amplitude}& $\xi_{lin}>|\xi_f|$& Float amplitude (linearity)& 3&Nonlin. ineq. constraint $g_{NL,22}$\\
  \requirement{slam_amplitude}& $\xi_{\text{slam}}>|\xi_f|$& Float amplitude (slamming)& 3&Nonlin. ineq. constraint $g_{NL,23}-g_{NL,239}$\\
  \requirement{float_hydrostatics}&$V_{f,\text{struct},\text{min}} < V_{f,\text{struct}} < V_{f,\text{struct},\text{max}}$& Float hydrostatics& 3&Nonlin. ineq. constraints $g_{NL,1}-g_{NL,2}$\\
  \requirement{spar_hydrostatics}&$V_{s,\text{struct},\text{min}} < V_{s,\text{struct}} < V_{s,\text{struct},\text{max}}$& Spar hydrostatics& 3&Nonlin. ineq. constraints $g_{NL,3}-g_{NL,4}$\\
  \requirement{FOS}&$FOS > FOS_{\text{min}}$& Structural factors of safety& 3&Nonlin. ineq. constraints $g_{NL,6}-g_{NL,13}$\\
  \requirement{GM}&$GM > 0$& Metacentric height (pitch stability)& 3&Nonlin. ineq. constraint $g_{NL,5}$\\
  \requirement{force_max}&$F_p = F_{\text{max}}$& Maximum powertrain force& 3&Nonlin. ineq. constraint $g_{NL,16}-g_{NL,17}$\\
  \requirement{P_max}&$P_{pk,\text{elec}} = P_{\text{max}}$& Maximum powertrain power& 3&Nonlin. ineq. constraint $g_{NL,18}$\\
  \requirement{LCOE}&$LCOE<LCOE_{\text{max}}$& Maximum LCOE& 3&Nonlin. ineq. constraint $g_{NL,15}$\\

    \caption{Requirements}
    \label{remdo:tab:requirements}
\end{longtable}
}
\ifdefined\DISSERTATION
    Priority 1 requirements are model requirements that must never be violated during the optimization, because doing so causes the simulation to error or return an invalid result such as \texttt{Inf}, \texttt{NaN}, or imaginary numbers.
    Some algorithms fail outright at such points, and even robust algorithms struggle to converge without a valid objective to guide the next step.
    Violating a priority 2 requirement yields an inaccurate but computationally valid objective; this is still undesirable during optimization (especially for gradient-based methods, where it can mislead the search), and an active priority 2 constraint in the final solution suggests the optimum lies in a region the model cannot accurately represent.
    Priority 3 requirements return an accurate objective even when violated, marking a design as undesirable for practical or functional reasons rather than model validity; active priority 3 constraints in the final design are acceptable and even expected.
    This three-level prioritization ultimately informs the selection of design variables and parameters.
\else
    Priority 1 requirements must never be violated during optimization, as they cause the simulation to error or return invalid results (\texttt{Inf}, \texttt{NaN}, imaginary numbers) that can stall or derail the algorithm.
    Priority 2 violations yield an inaccurate but valid objective, and constraints undesirable during the search and a sign, if active at the solution, that the optimum lies outside the model's accurate range.
    Priority 3 requirements return an accurate objective and constraints even when violated; active priority 3 constraints in the final design are acceptable and expected.
    This three-level prioritization informs the selection of design variables and parameters.
\fi

\ifdefined\DISSERTATION
    Requirements are ordered by priority in \Cref{remdo:tab:requirements}, and within priority grouped by enforcement mechanism.
    They include direct geometric relations among float, column, plate, and seafloor dimensions to prevent overlap (\reqref{float_spar_diam}, \reqref{float_in_out_diam}, \reqref{float_spar_draft}, \reqref{damping_plate_thickness}, \reqref{float_stiffener_height}), numerical overflow (\reqref{float_seafloor_overflow}, \reqref{spar_seafloor_overflow}), deviation from the truncated-cone shape (\reqref{float_base_diameter}, \reqref{float_draft}), submersion below the waterplane (\reqref{float_submergence}), and violation of dynamic assumptions (\reqref{spar_similarity}, \reqref{spar_damping_ratio});
    a power criterion ensuring net energy generation (\reqref{positive_net_power});
    height conditions for valid float-spar positioning in the static case (\reqref{clearance_static}, \reqref{tops_static}, \reqref{base_static}) and in operational seas (\reqref{tops_dynamic}, \reqref{base_dynamic}, \reqref{clearance_dynamic});
    maximum float amplitude based on linearity (\reqref{linear_amplitude}) and slamming (\reqref{slam_amplitude});
    buoyancy measures ensuring the float and spar float and can be ballasted (\reqref{float_hydrostatics}, \reqref{spar_hydrostatics});
    eight structural factors of safety (\reqref{FOS}); a stability measure preventing tipping (\reqref{GM});
    equality of maximum experienced and rated powertrain force and power to avoid oversized PTO components (\reqref{force_max}, \reqref{P_max});
    and a maximum LCOE constraint required for the epsilon-constraint procedure of \Cref{remdo:sec:multi-obj-process} (\reqref{LCOE}).
    Note that \reqref{base_static} dominates \reqref{float_spar_draft}, making the latter unnecessary to implement independently.
\else
    The requirements span geometric relations preventing overlap, numerical overflow, shape deviation, submersion, and dynamic-assumption violation (\reqref{float_spar_diam}--\reqref{spar_similarity});
    net power generation (\reqref{positive_net_power}); static and dynamic float-spar positioning (\reqref{clearance_static}--\reqref{clearance_dynamic});
    amplitude limits from linearity and slamming (\reqref{linear_amplitude}, \reqref{slam_amplitude});
    buoyancy and ballast feasibility (\reqref{float_hydrostatics}, \reqref{spar_hydrostatics});
    eight structural factors of safety (\reqref{FOS}); pitch stability (\reqref{GM});
    powertrain rating equality to avoid oversizing (\reqref{force_max}, \reqref{P_max});
    and a maximum-LCOE constraint for the epsilon-constraint procedure of \Cref{remdo:sec:multi-obj-process} (\reqref{LCOE}). 
    \reqref{base_static} dominates \reqref{float_spar_draft}, so the latter need not be implemented independently.
\fi
In summary, the requirements table clearly articulates and prioritizes the necessary items based on both design needs and model assumptions,
and will inform the optimization problem formulation.

%%%%%%%%%%%%%%%%%%%%%%%%
\subsubsection{Design Variables, Parameters, and Constraints}
The encoding of a physical design space into optimization design variables is non-unique.
\ifdefined\DISSERTATION
    For example, the float and spar drafts could be encoded directly as dimensions $T_{f,2}$ and $T_s$, or nondimensionally as dimension $T_s$ and ratio $T_{f,2}/T_s$.
    An encoding may be chosen because it more directly controls the objective, makes the problem convex, or helps maintain feasibility and convergence during optimization; the last is pursued here.
    Depending on the encoding, each requirement can be enforced via parameters, design variable bounds, linear inequality or equality constraints, or nonlinear inequality or equality constraints, listed from least to most difficult for the optimizer.
    It is therefore desirable to use only parameters, bounds, and linear constraints (no nonlinear constraints) for the five priority 1 requirements, ensuring they hold even at intermediate points.
    
    Continuing the draft example, two priority 1 requirements, \reqref{float_spar_draft} and \reqref{spar_seafloor_overflow}, contain $T_{f,2}$, while $T_s$ appears in \reqref{float_spar_draft} and \reqref{float_seafloor_overflow}.
    Assuming the other variables involved (e.g.\ $T_{f,1}$, $h$) are parameters or design variables, the dimensional encoding makes all three requirements achievable via linear constraint, while the nondimensional encoding makes \reqref{float_spar_draft} a bound (preferable) but forces \reqref{spar_seafloor_overflow} into a nonlinear constraint (undesired).
    Early results with nondimensional encodings \citep{mccabe_multidisciplinary_2022} showed the optimizer frequently violating nonlinear constraints and causing simulation failure, so the dimensional encoding is selected.
\else
    For example, bulk dimensions may be encoded directly as dimensions or as nondimensional ratios.
    An encoding may be chosen to control the objective directly, induce convexity in the objective or constraints, or maintain feasibility.
    The last is pursued here, since each requirement can be enforced, in increasing order of optimizer difficulty,
    via parameters, bounds, linear constraints, or nonlinear constraints,
    and it is desirable to keep all five priority 1 requirements free of nonlinear constraints so they hold even at intermediate points.
    Applied to the draft example, the dimensional encoding renders the relevant priority 1 requirements achievable by linear constraint whereas the nondimensional encoding forces a nonlinear constraint.
    Early results with nondimensional encodings \citep{mccabe_multidisciplinary_2022} showed frequent constraint violation and simulation failure, motivating the dimensional choice.
\fi

This logic repeats for the remaining design variables and requirements.
After several formulation iterations, the design variable and parameter definitions in \Cref{remdo:tab:design-vars,remdo:tab:parameters} emerged as effective.
The 12 design variables represent 5 bulk dimensions, 2 generator ratings, and 5 structural dimensions (3 thicknesses and 2 stiffener heights).
Referring back to \Cref{remdo:tab:lit}, this study is the most comprehensive WEC optimization in breadth of design-variable disciplines, and on the larger end, though not the largest, in number of design variables.

\begin{table}[ht]
\setlength\tabcolsep{1pt}
\renewcommand{\arraystretch}{1.1} % make taller
\begin{center}
{
\caption{Design Variables}\label{remdo:tab:design-vars}
\begin{tabular}{c>{\centering\arraybackslash}m{0.12\linewidth}>{\centering\arraybackslash}m{0.25\linewidth}>{\centering\arraybackslash}m{0.12\linewidth}>{\centering\arraybackslash}m{0.12\linewidth}>{\centering\arraybackslash}m{0.12\linewidth}c}
 \textbf{\#}&\textbf{Design Var.}& \textbf{Description} & \textbf{Lower Bound}& \textbf{Nominal Value} & \textbf{Upper Bound} & \textbf{Units} \\ \hline
 $x_1$&$D_s$& Spar outer diameter& 0& 6& 30& m\\
 $x_2$&$D_{f}$ & Float outer diameter& 1& 20 & 50& m \\
 $x_3$&$T_{f,2}$ & Float draft& 0.5& Nominal RM3~\citep{mccabe_development_2026}&100& m\\
 $x_4$&$h_s$& Spar height  & 5& Nominal RM3~\citep{mccabe_development_2026}& 100& m\\
 $x_5$&$h_{fs,\text{clear}}$& Vertical float tube to spar clearance at rest& 0.01& 4& 10& m\\
 $x_6$& $F_{\text{max}}$& Generator peak force rating & 0.01& 6.773 & 100& MN\\
 $x_{7}$&$P_{pk,\text{elec}}$& Generator peak power rating & 50& 286& 1000& kW\\
 $x_{8}$&$t_{f,b}$& Float bottom plate thickness & 2.5& 14.2& 25.4& mm\\
 $x_{9}$&$t_{s,r}$& Spar radial column thickness & 5& 25.4& 50.8& mm\\
 $x_{10}$&$t_{d}$& Damping plate thickness& 5& 25.4& 50.8& mm\\
 $x_{11}$&$h_{\text{stiff},f}$& Float stiffener height& 0& 0.406& 2& m\\
 $x_{12}$& $h_{\text{stiff},d}$& Damping plate stiffener height& 0& 0.559& 2&m\\
\end{tabular}%
}
\end{center}
\end{table}

The parameters in \Cref{remdo:tab:parameters} assume operation in the seas of Humboldt Bay, CA, A-36 structural steel, and nominal RM3 values for the design factor of safety, fixed charge rate, number of devices, and float submergence ratio \citep{RM3}.
\ifdefined\DISSERTATION
    Three parameters of note are the geometric ratios $D_d/D_s$, $T_s/D_s$, and $h_d/D_s$, which hold the spar proportions constant to maintain a sufficiently large damping ratio (\reqref{spar_damping_ratio}).
    Together with \reqref{spar_similarity}, which enforces geometric similarity on the damping plate, this fixes the spar dynamics to those of the nominal RM3, preserving the modeling assumptions of \citet{mccabe_development_2026}.
    In effect, the spar dimensional design space is restricted to comply with the model assumptions.
    This prevents determining the true optimal spar dimensions under multibody dynamics but suffices for an initial study emphasizing float design.
\else
    Three parameters of note, the geometric ratios $D_d/D_s$, $T_s/D_s$, and $h_d/D_s$, hold the spar proportions constant to maintain the damping ratio (\reqref{spar_damping_ratio}) and, with \reqref{spar_similarity}, fix the spar dynamics to those of the nominal RM3, preserving the modeling assumptions of \citet{mccabe_development_2026}.
    This restricts the spar design space to remain model-valid; it precludes finding the true optimal spar dimensions under multibody dynamics but suffices for an initial study emphasizing float design.
\fi
%\hl{Explain how this manifests as a minimum damping plate diameter.}

\begin{table}[ht]
\centering
\setlength\tabcolsep{2pt}
\renewcommand{\arraystretch}{1.4}
\caption{Selected Parameters}\label{remdo:tab:parameters}
\begin{tabular}{cP{0.28\linewidth}P{0.23\linewidth}c}
\textbf{Parameter} & \textbf{Description} & \textbf{Value} & \textbf{Units}  \\ \hline
$\rho_w$ & Seawater density& 1000 & kg/$m^3$  \\ 
$g$& Acceleration of gravity& 9.8 & m/$s^2$  \\ 
$h$& Water depth & 100 & m  \\ 
JPD& Wave joint probability distribution & read from file \citep{janzou_sam_2022}  &\%  \\ 
$H_s$ & Wave height \citep{janzou_sam_2022}& 0.25 : 0.5 : 6.75  & m  \\ 
$H_{s,\text{struct}}$ & 100 year wave height \citep{berg_extreme_2011}& [13.4, 18.8, 24.2, 30.1, 24.2, 18.8, 13.4] & m  \\ 
$T$& Wave energy period \citep{janzou_sam_2022}& 4.5 : 1 : 18.5  & s  \\ 
$T_{\text{struct}}$ & 100 year wave peak period \citep{berg_extreme_2011}& [5.57 8.76 12.18 17.26 21.09 24.92 31.70] & s  \\ 
$\sigma_y$ & Material yield strength& 248 & MPa  \\ 
$\rho_m$ & Material density& 7850 & kg/$m^3$  \\ 
$E$& Material Young's modulus& 200 & GPa  \\ 
$\text{cost}_m$ & Material cost& 1.89 & \$/kg \\ 
%$t_{ft}$ & \begin{tabular}[c]{@{}c@{}}Surface Float Top Plate\vspace{-1.5 mm} \\ Thickness\end{tabular} & 0.013 & m  \\ 
%$t_{fr}$ & \begin{tabular}[c]{@{}c@{}}Float Radial Wall\vspace{-1.5 mm} \\ Thickness\end{tabular} & 0.011 & m  \\ 
%$t_{fc}$ & \begin{tabular}[c]{@{}c@{}} Float Circumferential\vspace{-1.5 mm} \\ Gusset Thickness\end{tabular} & 0.011 & m  \\ 
%$t_{fb}$ & Float Bottom Thickness & 0.014 & m  \\ 
%$t_{sr}$ & Vertical Column Thickness & 0.025 & m  \\ 
%$B_{\text{min}}$ & Minimum Buoyancy Ratio & 1 & -  \\ 
$FOS_{\text{min}}$ & Minimum factor of safety& 1.5 & -  \\ 
$D_{d_{\text{min}}}$ & Minimum damping plate diameter & 30 & m  \\ 
$FCR$ & Fixed charge rate& 10.8 & \%  \\ 
$N_{WEC}$ & Number of WECs in array& 100 & -  \\ 
$D_{d}/{D_{s}}$ & Normalized damping plate diameter & 5 & -  \\ 
$T_{s}/{D_{s}}$ & Normalized spar draft& 5.83 & -  \\ 
$h_{d}/{D_{s}}$ & Normalized damping plate thickness & 0.004 & -  \\ 
$T_{f}/{h_{f}}$ & Float submergence ratio & 0.5 & -  \\ 
$F_{\text{heave},\text{mult}}$ & Heave force multiplier & 1.65 & - \\ 
\end{tabular}
\end{table}

Returning to \Cref{remdo:tab:requirements}, the enforcement column shows that this encoding successfully avoids nonlinear constraints for all priority 1 requirements and all but one priority 2 requirement, helping prevent model errors while facilitating convergence.
A total of \resultsRE[numLinConstraints]~linear and \resultsRE[numNonlinConstraints]~nonlinear inequality constraints exist, with no equality constraints; the eight linear inequality constraints are given in matrix form in \Cref{remdo:sec:appendix-linear-constraints}.

\ifdefined\DISSERTATION
    Notably, while the design variables include PTO sizing variables, PTO characteristics such as the generator and drivetrain impedances are kept as parameters.
    Rather than perform PTO co-design, which \citet{coe_useful_2023,gaebele_incorporating_2023} and studies discussed in \Cref{remdo:sec:lit}) have already proven is valuable, this study keeps the PTO design fixed at representative values.
    While impedance matching can and should be used to design the actual PTO hardware, this choice adds simplicity and illustrative value at the system level: it highlights the novel non-PTO (e.g.\ structural) aspects of WEC co-design introduced here while still capturing the effect of PTO rating on system performance and maintaining the generality to incorporate PTO optimization later.
    Developing a more detailed PTO cost model that scales with these characteristics would be a necessary next step toward PTO co-design.
\else
    While the design variables include PTO sizing (force and power ratings), other PTO characteristics such as generator and drivetrain impedances are kept as parameters rather than co-designed.
    This choice adds simplicity and isolates the novel non-PTO (e.g.\ structural) aspects of co-design introduced here, while still capturing the effect of PTO rating on system performance. 
    Recent papers have demonstrated the value of full PTO co-design \citep{coe_useful_2023,gaebele_incorporating_2023};
    pursuing it here would require a more detailed PTO cost model scaling with impedance, a necessary next step for future work.
\fi

%\subsubsection{Model Checking}
%\hl{todo: check for redundant constraints, monotonicity, etc} \citep{papalambros_principles_2017}.
%It was initially speculated that the structural constraints would always be active, 

\subsection{Interpretation}
Translating raw optimization outputs into meaningful engineering insights is nontrivial and rarely made explicit.
Consistent with MDO as part of the design process, this section outlines the methodology for interpreting optimation and sensitivity results for wave energy, focusing on which design questions can and cannot be addressed by which results.
% maybe instroduce this distinction when discussing choice of DVs vs parameters originally? think of PTO size, structures, and h_fs_clear as dynamic limit setters. maybe cite Pat's book on how parameter sensitivities can be used for factor screening and other things.
We separate parameters by the amount of control the designer has over them and the degree to which their values may vary due to both flexibility and uncertainty.
``States of the world,'' such as location and use case, have low control and high variation, whereas physical constants, such as gravitational acceleration and water density, have low control and low variation.
Performance-based characteristics are characterized by moderate control (through subsystem R\&D or market development) and can have high (i.e. cost uncertainty of unmodeled subsystems) or low variation (i.e. drag coefficient of a floating body).
Finally, parameters that could have been design variables but were instead fixed for simplicity, such as the damping plate diameter and spar draft, have high control and low variation.
